\documentclass[12pt,a4paper]{article}
\usepackage[UTF8]{ctex}
\usepackage[top=2.5cm,bottom=2.5cm,left=2.5cm,right=2.5cm]{geometry}
\usepackage{longtable}
\usepackage{array}
\usepackage{booktabs}
\usepackage{enumitem}

\renewcommand{\arraystretch}{1.5}

\pagestyle{empty}

\begin{document}

\begin{center}
{\Large\bfseries 计算机网络信息中心研究生学位论文修改说明}\\[6pt]
{\large （\CheckedBox 评阅修改意见\quad $\square$答辩修改意见）}
\end{center}

\vspace{6pt}

\noindent
\begin{tabular}{p{1.2cm}p{3cm}p{1.2cm}p{3cm}p{1.8cm}p{3cm}}
姓名 &  & 专业 &  & 联系电话 &  \\
类别 & \CheckedBox 硕士\quad $\square$博士 & 导师 &  &  &  \\
\end{tabular}

\noindent 论文题目：基于多元时序的分布式集群异常检测研究

\vspace{10pt}

%% ===== 专家一 =====
\noindent{\bfseries\large 专家一}
\vspace{4pt}

\begin{longtable}{|p{0.42\textwidth}|p{0.52\textwidth}|}
\hline
\centering\bfseries 修改意见 & \centering\arraybackslash\bfseries 修改情况 \\
\hline
\endfirsthead
\hline
\centering\bfseries 修改意见 & \centering\arraybackslash\bfseries 修改情况 \\
\hline
\endhead

1.在论文格式方面，个别章节存在不必要的空格，排版规范性有待进一步完善。
&
感谢专家的指正。已对全文进行了逐章排查，删除了各章节中不必要的空格和格式异常，统一了标点符号的使用规范（如中英文标点混用、公式前后多余空格等），修正了个别段落的缩进和对齐问题，排版规范性已进一步完善。
\\
\hline

2.在内容组织上，第五章作为系统实现部分，尚缺少整体系统框架图，对系统结构及各模块之间关系的展现不够直观，同时未能清晰说明前文所提出的算法如何在系统中进行集成与支撑，章节之间的衔接有待进一步加强。
&
感谢专家的宝贵建议。已从以下三个方面进行了修改：

（1）增加系统总述段落。在第5.1节开头新增了系统设计总述段落，从整体到局部地概述了系统的四层分层架构（用户层、应用层、算法层、数据层）、五个核心功能模块的职责边界及其协作关系，并明确说明了前文算法在系统中的集成方式：检测引擎模块作为算法与系统的桥梁，将第三章CSAD-AT算法和第四章NSFusion框架封装为可配置的检测流水线，根据检测场景自动编排对应的算法组件。

（2）新增系统总体架构图。新增了图5-1（集群节点异常检测系统总体架构），以分层结构图的形式直观展示了系统的四层架构、五个功能模块的位置关系和数据流向，以及第三章CSAD-AT和第四章NSFusion算法组件在算法层中的具体布局。

（3）强化章节衔接。扩写了第五章的引言段落，显式回顾了第三章和第四章的算法贡献及其在系统中的角色定位，将"算法研究——系统集成——工程落地"的逻辑链条表述得更加清晰。同时在5.1.1节功能模块设计和本章小结中均增加了对架构图的引用和对算法集成关系的呼应说明。
\\
\hline

\end{longtable}

\vspace{10pt}

%% ===== 专家二 =====
\noindent{\bfseries\large 专家二}
\vspace{4pt}

\begin{longtable}{|p{0.42\textwidth}|p{0.52\textwidth}|}
\hline
\centering\bfseries 修改意见 & \centering\arraybackslash\bfseries 修改情况 \\
\hline
\endfirsthead
\hline
\centering\bfseries 修改意见 & \centering\arraybackslash\bfseries 修改情况 \\
\hline
\endhead

1.3.4.2小节中分别从数值距离、嵌入空间距离和局部密度三个互补视角对节点的偏离程度进行量化，"其中密度偏离评分方面，与3.3.3节介绍的DBSCAN聚类方法不同，CSAD-AT算法借鉴局部离群因子的思想，设计了基于k近邻可达密度的连续异常评分方法"，鉴于此，又为何在3.3.3节提出基于DBSCAN聚类的密度相似性度量方法呢？请完善说明。
&
感谢专家的指正。3.3节与3.4节的关系是从基础方法介绍到算法设计的递进。3.3节分别从距离、嵌入和密度三个视角探索了独立可用的相似性度量方法，其中3.3.3节的DBSCAN方法能够自适应地将节点划分为正常簇和离群点，验证了密度分析视角在集群异常检测中的有效性。然而，DBSCAN输出的是离散的二值标签（簇内/噪声），无法产生连续的异常分数，不适合在3.4节的多视角融合框架中与其他视角的连续评分进行归一化聚合。因此，CSAD-AT算法在密度视角改用了基于k近邻可达密度的连续评分方法，使三个视角均输出连续分数，从而支撑后续的归一化加权聚合与I-SPOT自适应阈值机制。

原文对这一演进关系的交代不够充分，已在3.4.2节密度偏离评分的开头补充了从DBSCAN到连续评分方法的设计动机说明，明确指出DBSCAN的离散二值输出无法满足多视角融合框架对连续分数的要求，从而使3.3.3节与3.4.2节之间的逻辑衔接更加清晰。
\\
\hline

2.论文第4.5章节，请适当补充与其他同类型工作的对比。
&
感谢专家的宝贵意见。原稿实验部分仅选取了马氏距离作为基线方法，对比维度较为单一。根据专家建议，修改稿在4.5节中补充了条件变分自编码器（CVAE）作为深度生成模型类基线。CVAE通过将指标名称作为条件变量，在统一框架内学习全部监控指标的正常数据分布，以重建误差作为异常评分依据，代表了基于深度学习的异常检测技术路线。修改后的实验对比涵盖了经典多元统计方法（马氏距离）、深度生成模型（CVAE）以及所提方法的消融实验（单一数值方法、单一形状方法、双视角融合方法），从多个角度验证了NSFusion的有效性。

关于进一步扩展基线范围，本章所研究的问题具有特定的建模方式：其目标是在同一集群内依据多指标时间序列对各节点的异常程度进行排序，属于多节点横向对比的异常定位任务，与时序异常检测领域中常见的单实体异常片段检测设定存在本质区别。目前主流的开源时序异常检测算法（如Isolation Forest、LSTM-AE、OmniAnomaly等）主要面向单条或少量时间序列的异常区间识别，将其适配至本文的多节点排序场景需要对算法进行较大幅度的改造，改造后的方法难以代表原算法的设计意图，对比的公平性和参考价值有限。因此，修改稿选取了与本文问题建模最为契合的两类代表性基线进行实验对比。

修改内容见论文第4.5.1节第3段、第4.5.2节和第4.5.3节末段。
\\
\hline

3.论文第5.1小节，建议增加一个系统设计的总述段落以及系统架构图（包含5.1.1中介绍的各功能模块）。
&
感谢专家的建议。已在第5.1节开头新增了系统设计总述段落，从整体到局部地概述了系统的四层分层架构和五个核心功能模块的职责与协作关系，并明确说明了前文算法的集成方式。同时新增了图5-1（集群节点异常检测系统总体架构），以分层结构图直观展示了系统的四层架构、五个功能模块的位置关系和数据流向，以及CSAD-AT和NSFusion算法组件在算法层中的具体布局。（具体修改内容同专家一意见2的修改说明。）
\\
\hline

4.论文存在写作规范问题，如部分图中英文混杂、参考文献引用格式不规范、部分表述不够严谨、部分段落表述AI痕迹明显（如1.1小节最后一段内容）等，请修改完善。
&
感谢专家的细致审读。已针对所列各项问题逐一修改：（1）图中英文混杂：已将图2-5等图片中的英文标注替换为中文，全文图片标注已统一规范。（2）参考文献引用格式：已逐条检查并修正了作者姓名大小写、期刊名称缩写、卷号页码等格式问题，统一了引用样式。（3）AI痕迹修正：全文中"具体而言""值得注意的是"等程式化过渡词已替换或删除；第1.2节对仗式评述开头已改为直接阐述；第1.1节末段研究意义已整体重写，去除"创新性地""系统性地"等典型AI表述。修改涉及第一章、第三章和第四章。
\\
\hline

\end{longtable}

\vspace{10pt}

%% ===== 专家三 =====
\noindent{\bfseries\large 专家三}
\vspace{4pt}

\begin{longtable}{|p{0.42\textwidth}|p{0.52\textwidth}|}
\hline
\centering\bfseries 修改意见 & \centering\arraybackslash\bfseries 修改情况 \\
\hline
\endfirsthead
\hline
\centering\bfseries 修改意见 & \centering\arraybackslash\bfseries 修改情况 \\
\hline
\endhead

1.题目是基于多元时序的分布式集群异常检测研究，第三章和第四章题目中的多节点是与题目多元呼应么？但是并没有体现出来时序的概念，建议斟酌修改。
&
感谢专家的细致审阅与宝贵建议。

关于"多节点"与"多元"的关系：本文研究的核心数据形式是分布式集群中同一指标（或多个指标）在多个节点上同步采集的监控数据，每个节点的同一指标构成一条独立的时间序列，多个节点的同一指标共同构成一组多元时间序列。因此，"多节点"正是本文语境下"多元"的具体来源和物理含义——第三章的"单指标多节点"对应多元时序的受限情形（多元性来源于节点维度），第四章的"多指标多节点"则对应完整的多元时序场景（多元性同时来源于指标维度和节点维度）。使用"多节点"而非直接使用"多元"，是为了更准确地体现本文方法区别于传统多元时序异常检测的独特视角。

关于"时序"概念的体现：专家指出的问题十分准确。原标题确实未能在章标题层面体现"时序"这一核心概念。现已将第三章标题修改为"CSAD-AT：阈值鲁棒的单指标多节点时序异常检测算法"，第四章标题修改为"NSFusion：数值-形状融合的多指标多节点时序异常检测"，使章标题与论文大题目中"多元时序"的表述形成明确呼应。
\\
\hline

2.大模型生成语言太明显，例如，研究背景及意义第二段"运维人员常通过对比分布式集群中的同一指标曲线来判断异常：与其他变化模式相差很大的'离群曲线'更可能是异常"；第三段"具体而言，XX"，全文统一修改。
&
感谢专家的细致审读。已对全文进行了逐段排查，针对AI写作痕迹进行了系统性修正，主要包括以下几个方面：

（1）程式化过渡词的清理：将"具体而言""值得注意的是""总体来看""综上所述"等模式化连接词逐一替换为更自然的表述或直接删除，涉及第一章（2处）、第三章（3处）和第四章（3处），共计8处修改。例如，将"具体而言，本文从欧氏距离、自编码器和密度偏离三种视角..."改为"本文分别从欧氏距离、自编码器和密度偏离三种视角..."，去掉不必要的引导词。

（2）对仗式评述句式的改写：第1.2节研究现状综述中，"奠定了...理论基石""拓展了...实用边界"等工整对仗的总结性评述已改为直接、平实的学术阐述，避免人为修辞痕迹。例如，将"为时序异常检测奠定了坚实的理论基石"改为"推动了时序异常检测方法的发展"。

（3）专家标注段落的重写：第1.1节第二段中被专家标注的"运维人员常通过对比...来判断异常"一句，原文行文过于概括且带有AI生成的归纳式口吻，已重写为结合实际运维场景的具体描述；第三段中"具体而言"引出的内容已调整句式结构，直接切入方法描述。第1.1节末段研究意义部分也进行了整体重写，以更贴合论文实际贡献的表述替代原文中过度拔高的AI典型用语。

（4）全文其他零散修正：检查并修正了第二章文献综述中若干"近年来...取得了显著进展"等套话表述，以及第六章总结中"本文的主要贡献如下"等列举式开头，使全文语言风格更加统一自然。
\\
\hline

3.有些段落比较小，比如，3.3节上边一段话"上述结果表明，以学习单节点历史模式为核心思路的现有方法难以有效应对负载均衡类指标波动性强的挑战。"单独成段，建议与其他相关内容合并。
&
感谢专家的细致审阅。已对全文进行了逐段检查，将过短的独立段落与上下文相关内容进行了合并，具体修改如下：

（1）第3.2节：专家所指出的"上述结果表明，以学习单节点历史模式为核心思路的现有方法难以有效应对负载均衡类指标波动性强的挑战"原作为独立段落，现已将其并入前一段对基线方法实验结果的原因分析中，使问题总结与原因分析形成连贯的论述。

（2）第3.3节：CSAD-AT算法整体流程介绍部分存在两处过短段落，分别为评分融合步骤的简要说明和I-SPOT阈值判定的引出语，现已与各自对应的方法描述段落合并，保证方法介绍的叙述连贯性。

（3）第3.5节：消融实验部分"为验证各组件的贡献..."引出语原独立成段，现已并入后续实验设计的描述中；实验结论部分的单句总结段也已并入结果分析段落。

（4）第五章：系统实现部分原有多处功能模块的简短介绍各自独立成段（如数据采集模块、告警通知模块等），篇幅仅一两句话，现已按照功能层次将相关模块介绍合并为完整段落，使系统各层的功能描述更加充实紧凑。

修改后全文已无单句独立成段的情况，段落结构更加完整规范。
\\
\hline

4.全文图需要统一，包括字体大小一致、是否加粗一致，现在字体大小不和加粗都不一致。
&
感谢专家的指正。已对全文图片的字体规范进行了统一处理。重新检查并调整了各章图片中的文字标注，统一了字体大小和加粗样式，确保全文图片在视觉风格上保持一致。
\\
\hline

5.第三章公式后边如果没有加标点，则表示这句话还没说完，那么公式下边紧挨着的行不需要首行缩进，要顶格写。
&
感谢专家的指正。已对第三章全部12处公式环境逐一进行了检查和规范化处理。其中8处公式后无标点且后续文本确为同一句子的延续（以"其中""即"等承接），已确认格式正确（无空行、顶格书写）。另有4处公式（式3-1、式3-2、式3-3、式3-7）由冒号引出，公式本身即为句子结尾，但原文缺少句末标点，现已在公式末尾补充句号，以明确标示句子的完结。修改后全文公式标点与后续文本的缩进关系已统一规范。
\\
\hline

6.摘要中第一段提到"提出了基于多元时序曲线相似性的分布式"其中曲线相似性题目中并未提到，注意全文统一说法。
&
感谢专家的细致审阅。已接受该意见。原摘要第一段中"基于多元时序曲线相似性的"与论文题目的"基于"结构不一致，现已修改为"围绕多元时序的分布式集群异常检测展开研究"，与题目保持一致，"曲线相似性"改为在后文作为具体检测手段自然引出。英文摘要已同步修改。
\\
\hline

\end{longtable}

\vspace{20pt}

\noindent 导师审核意见：

\vspace{30pt}

\noindent 签字：\hspace{8cm}日期：

\vspace{10pt}

\noindent\small 说明：1.学生应对论文评阅/答辩意见做一一对应说明，表格可根据评阅专家/答辩委员会成员人数以及评阅/答辩意见的数量自行调整；\\
2.评阅意见修改、答辩意见修改分别提交。

\end{document}
